\section{Introduction}

Large language model systems for education are still commonly framed as retrieval-augmented generation pipelines: a learner question is matched to course materials or reference documents, relevant passages are retrieved, and a model produces an answer grounded in those passages. That pattern is effective when the task is primarily factual, citation-sensitive, or directly answerable from a curated corpus. Many educational tasks, however, are not merely information-access problems. Tutoring, misconception diagnosis, rubric-based critique, study planning, and adaptive scaffolding often require decomposition, role specialization, iterative critique, and judgments about how much help to provide at a given moment.

This paper therefore asks a narrower and more defensible question than whether multi-agent systems replace retrieval entirely. The central issue is whether multi-agent systems can replace or extend \emph{classical RAG as the top-level orchestration pattern} for educational tasks whose main difficulty lies in pedagogical coordination rather than document access alone. In this framing, retrieval remains important for evidence-grounded tasks, but it may no longer be sufficient as the primary architectural abstraction when tasks require explicit learner-state tracking, coordination among pedagogical roles, or multi-step instructional planning.

The manuscript is organized around two dimensions: the degree of knowledge grounding required by a task and the degree of pedagogical coordination complexity required to solve it well. The core hypothesis is conditional rather than universal. Classical RAG should remain strong on relatively simple knowledge-grounded tasks, while agentic or multi-agent systems become more compelling as tasks demand adaptive sequencing, critique, or role-specialized reasoning. This perspective avoids blanket replacement claims and instead emphasizes boundary conditions, mechanism, and fair architectural comparison.

The article makes three contributions. First, it proposes a task taxonomy that separates retrieval-heavy educational tasks from coordination-heavy ones. Second, it outlines an evaluation framework for comparing classical RAG, agentic RAG, and non-RAG multi-agent systems under shared budgets, source corpora, and task regimes. Third, it supplements that framework with repository-grounded evidence from executable tutoring architectures evaluated on EduBench and TutorEval snapshots, allowing the paper to move beyond a purely conceptual blueprint while still keeping claims disciplined.

\subsection{Research Questions}

The paper is organized around two research questions.

\begin{enumerate}[label=RQ\arabic*:,leftmargin=*]
    \item Under what educational task conditions is classical RAG sufficient, and when do agentic or multi-agent systems become justified?
    \item What evaluation design can compare classical RAG, agentic RAG, and non-RAG multi-agent systems fairly under shared tasks, budgets, and source corpora?
\end{enumerate}

\subsection{Paper Roadmap}

The remainder of the paper proceeds as follows. Section~2 reviews the relevant literature on educational evaluation, RAG, and multi-agent tutoring. Section~3 defines the task taxonomy and architectural comparison framework. Section~4 presents the evaluation design and benchmark mapping. Section~5 reports repository-grounded experimental snapshots on EduBench and TutorEval. Section~6 analyzes the resulting boundary conditions and ablation logic. Section~7 discusses what the current evidence does and does not establish, and Section~8 concludes with implications for future educational AI system design.
